% GENERATED FILE. Edit canonical lessons, then run npm run generate:books.
% canonical-source-hash: fnv1a64:100528f99f7365cd
% canonical-lessons: TE-C443-vahanam, TE-C443-injanu, TE-C443-tairu, TE-C443-breku, TE-C443-haranu

\chapter{Wheels and Engines}
\label{ch:te-wheels-and-engines}
\hlchaptermodality{443}

\begin{chapteropening}
\textbf{By the end of this chapter:} \emph{I can say a vehicle, an engine, a tyre, a brake and a horn.}

Complete the last lesson of chapter 443: I can say a vehicle, an engine, a tyre, a brake and a horn.
\end{chapteropening}

\section[\texorpdfstring{\te{వాహనం} (vāhanaṁ)}{వాహనం (vāhanaṁ)}]{\te{వాహనం} (vāhanaṁ) — a vehicle}

\label{lesson:TE-C443-vahanam}



\begin{quote}
Before the new one: say the Telugu for to be packed, then the Telugu for to derail.
\end{quote}

\subsection*{You'll want to know: \te{వాహనం}}
\textbf{\te{వాహనం}} — \emph{vāhanaṁ} — \textquotedblleft{}a vehicle\textquotedblright{}.

\begin{quote}
\te{ఆ} \te{రోడ్డు} \te{మీద} \te{పెద్ద} \te{వాహనాలు} \te{వెళ్ళకూడదు}.
\end{quote}

\begin{quote}
\emph{ā rōḍḍu mīda pedda vāhanālu veḷḷakūḍadu.}
\end{quote}

\begin{quote}
Big vehicles must not go on that road.
\end{quote}

Plural: \te{వాహనాలు}, \emph{vāhanālu}.

Plural \te{వాహనాలు}, the -\te{ం} noun plural in -\te{ాలు}. It covers every vehicle, from a bicycle to a lorry.

\begin{cousinweb}
\textbf{\te{వాహనం}} (\emph{vāhanaṁ}) is borrowed from Sanskrit \emph{vāhana}, \textquotedblleft{}any vehicle or draught animal\textquotedblright{}, from the root \emph{vah}, \textquotedblleft{}to carry\textquotedblright{} (CDIAL 11610).
\end{cousinweb}

\begin{grammarlens}[title={What you've built}]
One more word for getting around.
\end{grammarlens}

\subsection*{Your turn}
\begin{itemize}
\raggedright
  \item \emph{Say it:} \emph{vāhanaṁ}
  \item \emph{Say it:} \emph{vāhanaṁ}, once more
  \item \emph{Say it:} \emph{ā rōḍḍu mīda pedda vāhanālu veḷḷakūḍadu.} — Big vehicles must not go on that road.
  \item \emph{Recall:} say the Telugu for to be packed, then the Telugu for to derail, then say \emph{vāhanaṁ} again
\end{itemize}

\begin{tcolorbox}[breakable,colback=teal!4,colframe=teal!35!black,title={Before you move on}]
What does \te{వాహనం} mean? (\textquotedblleft{}A vehicle\textquotedblright{}.) Say it once more.
\end{tcolorbox}
\section[\texorpdfstring{\te{ఇంజను} (iñjanu)}{ఇంజను (iñjanu)}]{\te{ఇంజను} (iñjanu) — an engine}

\label{lesson:TE-C443-injanu}



\begin{quote}
Before the new one: say the Telugu for to derail, then the Telugu for a vehicle.
\end{quote}

\subsection*{You'll want to know: \te{ఇంజను}}
\textbf{\te{ఇంజను}} — \emph{iñjanu} — \textquotedblleft{}an engine\textquotedblright{}.

\begin{quote}
\te{బస్సు} \te{ఇంజను} \te{చాలా} \te{శబ్దం} \te{చేస్తోంది}.
\end{quote}

\begin{quote}
\emph{bassu iñjanu cālā śabdaṁ cēstōndi.}
\end{quote}

\begin{quote}
The bus engine is making a lot of noise.
\end{quote}

The engine of a train is also its \te{ఇంజను}.

\begin{cousinweb}
\textbf{\te{ఇంజను}} (\emph{iñjanu}) is a loan from English \emph{engine}, with the Telugu final -\te{ఉ}.
\end{cousinweb}

\begin{grammarlens}[title={What you've built}]
One more word for getting around.
\end{grammarlens}

\subsection*{Your turn}
\begin{itemize}
\raggedright
  \item \emph{Say it:} \emph{iñjanu}
  \item \emph{Say it:} \emph{iñjanu}, once more
  \item \emph{Say it:} \emph{bassu iñjanu cālā śabdaṁ cēstōndi.} — The bus engine is making a lot of noise.
  \item \emph{Recall:} say the Telugu for to derail, then the Telugu for a vehicle, then say \emph{iñjanu} again
\end{itemize}

\begin{tcolorbox}[breakable,colback=teal!4,colframe=teal!35!black,title={Before you move on}]
What does \te{ఇంజను} mean? (\textquotedblleft{}An engine\textquotedblright{}.) Say it once more.
\end{tcolorbox}
\section[\texorpdfstring{\te{టైరు} (ṭairu)}{టైరు (ṭairu)}]{\te{టైరు} (ṭairu) — a tyre}

\label{lesson:TE-C443-tairu}



\begin{quote}
Before the new one: say the Telugu for a vehicle, then the Telugu for an engine.
\end{quote}

\subsection*{You'll want to know: \te{టైరు}}
\textbf{\te{టైరు}} — \emph{ṭairu} — \textquotedblleft{}a tyre\textquotedblright{}.

\begin{quote}
\te{కారు} \te{టైరులో} \te{గాలి} \te{తక్కువగా} \te{ఉంది}.
\end{quote}

\begin{quote}
\emph{kāru ṭairulō gāli takkuvagā undi.}
\end{quote}

\begin{quote}
The car tyre is low on air.
\end{quote}

To pump air into it is \te{గాలి} \te{కొట్టు}.

\begin{cousinweb}
\textbf{\te{టైరు}} (\emph{ṭairu}) is a loan from English \emph{tyre}, with the Telugu final -\te{ఉ}.
\end{cousinweb}

\begin{grammarlens}[title={What you've built}]
One more word for getting around.
\end{grammarlens}

\subsection*{Your turn}
\begin{itemize}
\raggedright
  \item \emph{Say it:} \emph{ṭairu}
  \item \emph{Say it:} \emph{ṭairu}, once more
  \item \emph{Say it:} \emph{kāru ṭairulō gāli takkuvagā undi.} — The car tyre is low on air.
  \item \emph{Recall:} say the Telugu for a vehicle, then the Telugu for an engine, then say \emph{ṭairu} again
\end{itemize}

\begin{tcolorbox}[breakable,colback=teal!4,colframe=teal!35!black,title={Before you move on}]
What does \te{టైరు} mean? (\textquotedblleft{}A tyre\textquotedblright{}.) Say it once more.
\end{tcolorbox}
\section[\texorpdfstring{\te{బ్రేకు} (brēku)}{బ్రేకు (brēku)}]{\te{బ్రేకు} (brēku) — a brake}

\label{lesson:TE-C443-breku}



\begin{quote}
Before the new one: say the Telugu for an engine, then the Telugu for a tyre.
\end{quote}

\subsection*{You'll want to know: \te{బ్రేకు}}
\textbf{\te{బ్రేకు}} — \emph{brēku} — \textquotedblleft{}a brake\textquotedblright{}.

\begin{quote}
\te{డ్రైవరు} \te{వెంటనే} \te{బ్రేకు} \te{వేశాడు}.
\end{quote}

\begin{quote}
\emph{ḍraivaru veṇṭanē brēku vēśāḍu.}
\end{quote}

\begin{quote}
The driver braked at once.
\end{quote}

\te{బ్రేకు} \te{వేయు}, to put on the brakes; \te{బ్రేకులు} \te{పనిచేయలేదు}, the brakes failed.

\begin{cousinweb}
\textbf{\te{బ్రేకు}} (\emph{brēku}) is a loan from English \emph{brake}, with the Telugu final -\te{ఉ}.
\end{cousinweb}

\begin{grammarlens}[title={What you've built}]
One more word for getting around.
\end{grammarlens}

\subsection*{Your turn}
\begin{itemize}
\raggedright
  \item \emph{Say it:} \emph{brēku}
  \item \emph{Say it:} \emph{brēku}, once more
  \item \emph{Say it:} \emph{ḍraivaru veṇṭanē brēku vēśāḍu.} — The driver braked at once.
  \item \emph{Recall:} say the Telugu for an engine, then the Telugu for a tyre, then say \emph{brēku} again
\end{itemize}

\begin{tcolorbox}[breakable,colback=teal!4,colframe=teal!35!black,title={Before you move on}]
What does \te{బ్రేకు} mean? (\textquotedblleft{}A brake\textquotedblright{}.) Say it once more.
\end{tcolorbox}
\section[\texorpdfstring{\te{హారను} (hāranu)}{హారను (hāranu)}]{\te{హారను} (hāranu) — a horn}

\label{lesson:TE-C443-haranu}



\begin{quote}
Before the new one: say the Telugu for a tyre, then the Telugu for a brake.
\end{quote}

\subsection*{You'll want to know: \te{హారను}}
\textbf{\te{హారను}} — \emph{hāranu} — \textquotedblleft{}a horn\textquotedblright{}.

\begin{quote}
\te{వెనక} \te{ఉన్న} \te{కారు} \te{హారను} \te{కొట్టింది}.
\end{quote}

\begin{quote}
\emph{venaka unna kāru hāranu koṭṭindi.}
\end{quote}

\begin{quote}
The car behind sounded its horn.
\end{quote}

\te{హారను} \te{కొట్టు}, to sound the horn: the same \te{కొట్టు} as in to hit or to beat.

\begin{cousinweb}
\textbf{\te{హారను}} (\emph{hāranu}) is a loan from English \emph{horn}, the warning horn of a vehicle.
\end{cousinweb}

\begin{grammarlens}[title={What you've built}]
One more word for getting around.
\end{grammarlens}

\subsection*{Your turn}
\begin{itemize}
\raggedright
  \item \emph{Say it:} \emph{hāranu}
  \item \emph{Say it:} \emph{hāranu}, once more
  \item \emph{Say it:} \emph{venaka unna kāru hāranu koṭṭindi.} — The car behind sounded its horn.
  \item \emph{Recall:} say the Telugu for a tyre, then the Telugu for a brake, then say \emph{hāranu} again
\end{itemize}

\begin{tcolorbox}[breakable,colback=teal!4,colframe=teal!35!black,title={Before you move on}]
What does \te{హారను} mean? (\textquotedblleft{}A horn\textquotedblright{}.) Say it once more.
\end{tcolorbox}
